\section{Two pressures on the cylinder}\label{c4:pressure}

The finite marked identity will determine a positive winding-polygon mass
only after we know that its leading coefficient does not vanish. We obtain
an independent lower bound from a derivative of the cylinder partition
function, and use another parameter derivative to measure the unrooted planar mass
lost under projection to the cylinder.

Use the physical vectors \(h,v\) and cells \(P_{ij},R_{ij}\) of
Section~\ref{c4:cylinder}. For sufficiently large \(N\), on the cylinder
of period \(Nh\), count
unoriented polygons modulo axial translation by \(v\) only, with weight
\(w(\gamma)=x_*^{|\gamma|}\). Write \(/v\) for these axial classes and set
\[
 A_N=\sum_{\mathrm{winding}\ \gamma/v}w(\gamma),\qquad
 A_N^0=\sum_{\mathrm{contractible}\ \gamma/v}w(\gamma).
\]
These are initially nonnegative sums; the fixed-\(N\) argument below
proves their finiteness. Our first target is \(A_N\sim c/N\), with
\(c>0\). The finite-size correction to \(A_N^0/N\) will then give a
positive planar projection loss of order \(N^{-2}\). Both calculations
use the same fusion relation, but different boundary data. The final
subsection derives transverse-span bounds that turn these cylinder
estimates into planar tail estimates.

In this section set \(d=d_*=\pi/8\), \(q_\lambda=e^{2i\lambda}\); thus
\(q=q_d\). Write \(y=u^2\) for the squared auxiliary parameter, as before
\(t_i=z_i^2\) for squared spatial parameters. Angle letters and kernels
below are section-local notation.

\subsection{Finite responses and their physical meaning}

On the diagram matching space of the cylinder the closed transfer has the two loop weights \(n=-2\cos(4\lambda)\) and \(s=\eta+\eta^{-1}\). Use the analytic twist branch near \(\eta=i\) at \(s=0\). For each fixed \(N\), fixing the auxiliary initially at \(u=Q^{-1}\), the physical rank-one convergence on matchings proved in Section~\ref{c4:cylinder} gives a simple isolated eigenbranch near \(1\), analytic for sites near all \(1\), \(\lambda\) near \(d\), and \(s\) near zero, with diagram eigenvector normalized by empty coefficient \(1\). Its spin image is nonzero, and the corresponding eigenspace \emph{in the link image} is one-dimensional (it is a quotient representation for each fixed choice of parameters). By link-image invariance and transfer commutation, this spin vector is a common eigenvector for all auxiliaries. Call its eigenvalue \(\Lambda(y)\). This is rational in \(y\) by parity, with denominator dividing \(\prod_i(y-q_\lambda^2t_i)(y-q_\lambda^3t_i)\), bounded at \(0,\infty\); after that common denominator is used, the coefficients are analytic in the data under variation here. Near physical \(y\) it still gives the leading simple analytic diagram eigenvalue: the diagram eigenvalue there has nonzero spin image near the base point, and the branches agree by isolation on the quotient (or by joint commutation on spin images and separation from the other diagram eigenvalues). More explicitly, \(\Lambda(y)\) varies continuously near the value \(1\) there and as an image eigenvalue belongs to the diagram spectrum. At \(\lambda=d,s=0\) in fact \(\Lambda(y)=1\) identically, including at nearby lists of sites: the empty coefficient is fixed by transfer on the link image for every \(y\).

We use the twist responses
\(J_m(y)=[s^m]\Lambda(y)|_{\lambda=d}\), \(m\ge1\), and the
contractible response
\(J(y)=\partial_\lambda\Lambda(y)|_{\lambda=d,s=0}\).
At homogeneous sites their physical evaluations are
\[
 J_1(q^{-1})=A_N,\qquad
 J(q^{-1})=N(\log t_{\rm phys})'(d)+8A_N^0,
\]
where \(t_{\rm phys}(\lambda)\) is the empty-cell weight at
\(u=Q_\lambda^{-1}\).

To prove these identities, keep \(N\) fixed and place vacuum boundaries
on a finite-height cylinder. Its partition function is exactly the gas
of vertex-disjoint polygons with step weight \(X_\lambda\), loop weights
\(n,s\), and the empty scalar \(t_{\rm phys}\) per cell. The log
partition function per row has derivatives tending, as the height grows,
to those of the leading log eigenvalue. Indeed the vacuum matrix element
of its leading projection is analytic and nonzero nearby (it equals 1
at the base point), while the remaining powers divided by the leading
one decay exponentially on a small parameter disc, by the fixed-\(N\)
spectral gap and resolvent contours. At the base point \(\Lambda=1\), so
the first derivatives of \(\log\Lambda\) and \(\Lambda\) agree. The
motion of the physical spectral parameter in the \(\lambda\)-derivative
contributes zero because \(\Lambda(y)=1\) for every \(y\) there.
Differentiating at vanishing loop weights selects one polygon, with
\(n'(d)=8\); the empty cells contribute
\(N(\log t_{\rm phys})'(d)\). After division by height, each axial
orbit has a placement fraction increasing to 1. Monotone convergence of
these positive single-polygon sums proves both identities and their
finiteness. The subsequent analysis concerns their behavior as
\(N\to\infty\).

\subsection{Fusion constraints and convolution operators}

At \(y=\infty,0\) respectively,
\[
             \Lambda(y)=(1+s)(-e^{\pm 8i\lambda})^N.
\]
Indeed one can only use the limiting weighted swaps in the twisted trace (the other one-way flip of the auxiliary cannot close), then spatial total charge is zero. The two product identities give
\[
 \Lambda(t_j)\Lambda(q_\lambda^{-2}t_j)
    =\Big(\prod_i d_f(x_{ij})\Big)\Lambda(q_\lambda^{-1}t_j),\qquad
 \Lambda(t_j)\Lambda(q_\lambda^{-3}t_j)=\prod_i D(x_{ij})
\]
in the earlier tensor notation. All the sampled arguments here are away from poles near the base point. We need two types of responses:
\par\noindent\textbullet\quad For the twist responses, \(d_f,D\) are
\(1\), so \(\Lambda(qy)\Lambda(q^{-1}y)-\Lambda(y)\) has zeros at both
\(q^{-1}t_j,q^{-2}t_j\).
\par\noindent\textbullet\quad For the contractible response \(J\), the denominator needs no extra orders from pole motions since \(\Lambda-1=0\) identically at \(d\). Differentiating the products, argument motions in \(\Lambda\) likewise drop out. Now
\[
          J(qy)+J(q^{-1}y)-J(y)- S_X(y),\qquad
          S_X(y)=8i\sum_l (y-t_l)/(y+t_l)
\]
has the same zeros. Indeed \(\partial_\lambda d_f(x)|_d=8\cot(3d-x)\), and \(D(x)=d_f(x)d_f(x+\lambda)\). Writing \(z_l=e^{i\ell_l}, y=e^{2i\nu}\), the first product derivative gives \(8\sum_l\cot(3d-\ell_l+\ell_j)\) at \(\nu=\ell_j-d\). Subtract it from the second derivative to get \(8\sum_l\cot(2d-\ell_l+\ell_j)\) at \(\nu=\ell_j-2d\). Both are the indicated \(S_X\).

In particular at homogeneous sites these are order-\(N\) zeros at each of the two arguments (for the first assertion also coefficientwise), by separating the sites and analytic confluence.

Hereafter use homogeneous sites in this section, and put
\[
        y=e^{-3id}e^w,\qquad h_0=3d,\qquad \alpha=\pi/(6d)=4/3 .
\]
View the responses as functions of \(w\), so that the physical evaluation
\(y=q^{-1}\) is now \(w=id\). They have poles only at
\(w=\pm id+i\pi\) modulo \(2\pi i\), of order at most \(N\), and
endpoints as \(\Re w\to\pm\infty\) given by the extreme swaps. We first
give the convolution construction common to both responses. Its
accretivity will control their boundary norms; its supported factors will
later determine the signs of the two limiting constants.

Write \(H_0(w)=e^w/(1+e^w)\). Each such function \(V(w)\) (a \(J_m\) or
\(J\)) expands as a constant plus a sum of
\(\partial_w^k H_0(w\pm id)\), \(0\le k<N\), with scalar coefficients,
by ordinary rational partial fractions. Form \(v(w)\) from the expansion
by replacing \(H_0(w\pm id)\) by \(H_0(\alpha(w\pm id))\) under the
derivatives (same endpoints as \(V\)). Then on horizontal lines
\(|\Im w|\le h_0\)
\begin{equation}
 V(w+2id)+V(w-2id)-V(w)=\mathsf T*v,\qquad
 \widehat{\mathsf T}(k)=\alpha(2\cosh(2dk)-1)\frac{\sinh(6dk)}{\sinh(8dk)},\quad
 \widehat{\mathsf T}(0)=1,                                                 \label{c4:Press-1}
\end{equation}
where convolution is real translation, transform convention \(e^{-ikx}\). Indeed the transform of \(H_0'\) is \(\pi k/\sinh(\pi k)\) (e.g. shift a contour by \(2\pi i\) past the double pole), so this holds first on derivatives, with constants then fixed by endpoints. The symbol can also be written
\[
 \widehat{\mathsf T}(k)/\alpha
       =1-\big(\sinh(3\pi k/4)-\sinh(\pi k/2)\big)/\sinh(\pi k).
\]
Thus \(\mathsf T\) is a real even measure with point mass \(\alpha\) at 0 plus a kernel integrable with exponential weights of any rate \(<1\) (shift Fourier lines for the exponentially decaying analytic correction). In each norm
\[
 \|b\|_\pm^2=\sum_{\phi=\pm h_0}\int_{\mathbb R}|b(x+i\phi)|^2\alpha e^{\pm\alpha x}dx
\]
(the norm sign chosen independently of the summation), \(\mathsf T\) is bounded and strictly accretive: \(\Re\langle b,\mathsf T*b\rangle_\pm\ge c_T\|b\|_\pm^2\) for some \(c_T>0\). Indeed Plancherel after weighting uses symbol lines \(k\pm i\alpha/2\) for real \(k\). The fraction subtracted in the last display has modulus uniformly less than 1 there: squared denominator modulus minus squared numerator modulus is, with \(l=\pi k\),
\[
 \cosh^2 l-\cosh^2(3l/4)+\cosh(l/2)(\sqrt3\cosh(3l/4)-\cosh(l/2))>0,
\]
and the ratio decays at infinity.

We also use two convolution facts. First
\(\mathsf T=\mathsf T_+*\mathsf T_-\), with real measures supported
respectively on \([0,\infty),(-\infty,0]\), having inverses likewise
supported, all with exponential moments at some rate \(>\alpha/2\)
(in total variation). In fact the symbol is nonzero on \(|\Im k|<1\)
by \eqref{c4:Press-1}, positive on the real line. Its analytic logarithm,
real there, differs from \(\log\alpha\) by an exponentially decaying term
on smaller closed strips. Thus the logarithm is itself a transform of a
real measure with the required moments; split it by sign of position and
exponentiate in the convolution algebra. The resulting measures may be
signed. Their masses, however, are strictly positive, as are their real
Laplace moments whenever the corresponding exponential weight is bounded
on the support; these numbers are exponentials of real numbers in the
same construction. Second, writing
\(F=V(\cdot+2id)+V(\cdot-2id)-V\), one recovers
\[
       V=L*F,\qquad \widehat L(k)=(2\cosh(2dk)-1)^{-1}.
\]
This again follows first on derivatives and then by endpoints. The kernel shifts \(L(x+ip)\), \(|p|\le 2d-\epsilon\), are \(O_\epsilon(e^{-\alpha|x|})\); more precisely
\[
 L(\pm x+ip)=c_L e^{-\alpha x}e^{\mp i\alpha p}+O_\epsilon(e^{-2\alpha x}),\quad x\to+\infty,\qquad
 c_L=1/(2\sqrt3 d),
\]
by passing the first simple Fourier poles. One can shift \(L*F\) using lines of \(F\) interior to the strip of half-width \(h_0\) and shifting \(L\) by less than \(2d\) (boundedness on intermediate lines at fixed data, or analytic continuation).

\subsection{Essential twist response}\label{c4:essential-pressure}

The first response determines \(A_N\). Higher responses will control the
weight of tall winding polygons through the connected coefficients of the
finite gas. We first bound every fixed order, and then identify the
positive leading constant at first order.

Write \(f_m\) for \(v\) from \eqref{c4:Press-1} with \(V=J_m\), and set
\[
 F_m=\mathsf T*f_m,\qquad
 P_m=\sum_{0<j<m}J_j(w+2id)J_{m-j}(w-2id),\qquad G_m=F_m+P_m.
\]
The limits at either end are \(c_m=[m=1]\) for \(J_m,f_m,F_m\),
\(p_m=[m=2]\) for \(P_m\), and \(d_m=c_m+p_m\) for \(G_m\).
Fusion gives \(G_m\) the order-\(N\) zeros at \(w=\pm id\).
Use
\[
          z=e^{\alpha w},\qquad
          B_N(z)=\left[\frac{(z-a)(z-\bar a)}{(z+a)(z+\bar a)}\right]^N,\qquad a=e^{i\pi/6}.
\]
Then \(B_N\) is inner on the right half-plane and invariant by \(z\mapsto1/z\). The two norms are ordinary \(L^2\) on the imaginary boundary in \(z,1/z\) respectively. There is an orthogonality in either:
\[
                 \langle f_m-c_m B_N,\ G_m-d_m B_N\rangle_\pm=0 .
\]
Indeed \(f_m/B_N-c_m\) is rational analytic left, poles there canceled by scaling the partial fraction basis exactly as above, while \(G_m/B_N-d_m\) is analytic right by fusion (take the principal log). In either chosen coordinate both are bounded at zero and decay at infinity at least \(O(|z_{\rm coord}|^{-1/\alpha})\); thus divide both factors by the boundary phase \(B_N\), conjugate-reflect the left function and close a right semicircle. Also \(\|B_N-1\|_\pm=O(\sqrt N)\) by the boundary pointwise bound \(\min(2,CN\min(|z|,|z|^{-1}))\).

For fixed \(m\) we claim, inductively,
\begin{equation}
 \|f_m-c_m\|_\pm+\|G_m/B_N-d_m\|_\pm\le C_m\sqrt N,\qquad
 |J_m(x+i\phi)|\le C_{m,\epsilon}\min(1,e^{\alpha|x|}/N)\quad (|\phi|\le 5d-\epsilon).
                                                                            \label{c4:Press-2}
\end{equation}
Given the earlier bounds, \(\|P_m-p_m\|_\pm\le C_m\sqrt N\) (zero for \(m=1\)): of the two shifted arguments at each boundary one is safe pointwise by induction \((\phi=\pm d)\), and both have \(J_j-c_j\) norm \(O_j(\sqrt N)\) on their respective shifts. For the latter, transport \(f_j-c_j\) at the boundary to \(J_j-c_j\) at displacement \(ip\), \(|p|\le2d\), using the bounded weighted Fourier multiplier
\[
                e^{-p k}\alpha\sinh(6dk)/\sinh(8dk)
                 \quad\text{at }k\in\mathbb R\pm i\alpha/2.
\]
This follows from the same derivative transforms by contour shift (the endpoint-subtracted functions and derivatives decay faster than the weights require). Now \(G_m-d_m B_N\) equals \(\mathsf T*(f_m-c_m B_N)\) plus errors \(P_m-p_m, c_m(\mathsf T-I)*(B_N-1), -p_m(B_N-1)\). Accretivity and orthogonality yield the two norm bounds. By Cauchy interior evaluation of \(G_m/B_N-d_m\) in \(z,1/z\), the correction to \(d_m\) there is bounded by \(C_{m,\epsilon}\sqrt{N e^{-\alpha|x|}}\) for \(|\phi|\le h_0-\epsilon\). Restoring \(B_N\) suppresses this exponentially, giving
\[
           |G_m(x+i\phi)|\le C_{m,\epsilon}\exp(-c_\epsilon N e^{-\alpha|x|}) .
\]
As \(P_m\) there is bounded by \(C_{m,\epsilon}\min(1,e^{2\alpha|x|}/N^2)\), so is \(F_m\). Recover \(J_m=L*F_m\) with the indicated line shifts to reach \(5d-\epsilon\). Convolution with the exponential kernel bound gives \eqref{c4:Press-2} (for \(e^{\alpha|x|}<N\) split translations at absolute distance \(\alpha^{-1}\log N-|x|\)). This closes the induction. In particular every fixed positive order coefficient of \(\log\Lambda(q^{-1})\) is \(O_m(N^{-1})\).

\paragraph{The first response at the two edges.}
We now sharpen the first coefficient, for which \(F_1=G_1\). In each
edge coordinate \(u_e=N/z\) (right edge in \(x\)) or \(Nz\) (left edge),
put \(b_N(u_e)=B_1(u_e/N)^N\to b_\infty(u_e)=e^{-b u_e}\),
\(b=2\sqrt3\). Thus \(B_N(z)=b_N(u_e)\). Form
\[
             g_N=(f_1-b_N)/u_e,\qquad h_N=(F_1/b_N-1)/u_e.
\]
They are analytic in the right half-plane with uniformly bounded boundary
\(L^2(du_e)\) norms by \eqref{c4:Press-2} and the orthogonality estimates
after change of variables. Also \(g_N/b_N\) is analytic in the left
half-plane. For each right function its boundary values on \(u_e=i y'\)
have Fourier inverse supported on \([0,\infty)\), and interior values
are the Laplace transform of that density. Indeed one may close contours
(the functions are \(O_N(|u_e|^{-1-1/\alpha})\) at infinity and
\(O_N(|u_e|^{-1+1/\alpha})\) at zero), or use the Cauchy integral with
denominator \(v'-u_e\) for evaluation at \(v'\) on the right, which by
Plancherel picks out the nonnegative half-density. Similarly
\(g_N/b_N\) has density supported on the negative half-line.

Take weak \(L^2\) subsequences for \(g_N,h_N\). The support conditions
are weakly closed, and Cauchy evaluation is a bounded functional, so their
interior values converge. Division by the inner phases
\(b_N\to b_\infty\) converges strongly as a boundary multiplier.
Consequently the limit \(g\) still has nonnegative density support,
while \(g/b_\infty\) has nonpositive density support. Since division by
\(e^{-bu_e}\) shifts the density left by \(b\), the density of \(g\)
is supported in \([0,b]\).

To read the boundary data for the unscaled limits, use
\[
             \frac{b_\infty(u_e)}{u_e}
                    =\int_b^\infty e^{-u_e t}\,dt.
\]
The limits \(f_1\to b_\infty+u_e g\) and
\(F_1\to b_\infty(1+u_e\lim h_N)\) can therefore be written
\[
 f_{\rm edge}(u_e)=u_e\int_0^\infty e^{-u_e t}\varphi(t)dt,\qquad
 F_{\rm edge}(u_e)=u_e\int_0^\infty e^{-u_e t}\psi(t)dt ,
\]
with \(\varphi-1,\psi-1\in L^2(dt)\). The density of \(g\) is confined
to \([0,b]\), whereas multiplication by \(b_\infty\) shifts the density
of \(\lim h_N\) from \([0,\infty)\) to \([b,\infty)\). Hence, almost
everywhere,
\[
                 \varphi(t)=1\quad(t>b),\qquad
                 \psi(t)=0\quad(t<b).
\]

\paragraph{Determining the edge limit.}
Relation \eqref{c4:Press-1} passes to these limits on \(u_e>0\), since
\(f_1\) in the edge coordinate is bounded by \(1+C\sqrt{u_e}\) and
\(\mathsf T\) has enough exponential moments. The convolution dilates
\(u_e\) by \(e^{\alpha s'}\), with \(s'\) integrated against
\(\mathsf T(ds')\) (either edge by evenness). Laplace uniqueness gives
\(\psi=\mathsf T*\varphi\) in the coordinate
\(x'=\log(t/b)/\alpha\). These operations are bounded on constants plus
\(L^2(dt)\): the density dilation
\(\varphi(t)\mapsto\varphi(te^{-\alpha s'})\) has \(L^2(dt)\) norm
multiplier \(e^{\alpha s'/2}\), covered by the stated exponential moments.

The factorization now joins the two prescribed half-lines. In \(x'\), set
\[
       q_{\rm edge}=\mathsf T_+^{-1}*\psi=\mathsf T_-*\varphi.
\]
For \(x'<0\), the first expression vanishes by the supports of
\(\mathsf T_+^{-1}\) and \(\psi\). For \(x'>0\), every argument sampled
by the second expression lies where \(\varphi=1\), so it equals
\(c'=\int\mathsf T_-(ds')>0\). Thus
\(q_{\rm edge}=c'\mathbf1_{(0,\infty)}\), which determines the limit
uniquely. Applying \(\mathsf T_+\) and integrating gives
\[
 \int_0^\infty F_{\rm edge}(u_e)du_e=\int_b^\infty \psi(t)t^{-2}dt
       =\frac{c'}b\int e^{-\alpha s'}\mathsf T_+(ds')>0 .
\]
Absolute integrability follows from support and the stated \(L^2\)
bounds; positivity follows from the real Laplace moments of the factors.
Both edges have the same limit, without subsequence dependence. To return
to the winding activity, use \(J_1(id)=\int L(id-x)F_1(x)dx\).
Changing variables to \(u_e=N e^{-\alpha|x|}\) on the two halves,
the bound \(|F_1(x)|\le C\exp(-c u_e)\) and the shifted
\(L\)-asymptotic give by domination
\begin{equation}
           N J_1(id)\ \longrightarrow             2 c_L \alpha^{-1}\cos(\alpha d)\int_0^\infty F_{\rm edge}(u_e)du_e>0 .
                                                                             \label{c4:Press-3}
\end{equation}

\subsection{Contractible response}\label{c4:contractible-pressure}

For the second physical response we must separate its bulk term from the
finite-size correction. Use \(J=\partial_\lambda\Lambda|_d\) at
\(s=0\); its right/left endpoints are \(j_R=8iN,j_L=-8iN\). Let \(f\)
be its scaled partial fraction version, \(F=\mathsf T*f\), and
\[
      S(w)=8iN(2H_0(w-3id)-1),\qquad G=F-S,\qquad g=\mathsf T^{-1}*S.
\]
Here \(G\) has the order-\(N\) zeros, and \(f,g,F,S\) share the endpoints of \(J\). The forced term \(g\) is analytic throughout a slightly wider strip than \(|\Im w|\le h_0\), and
\[
 \widehat{g'}(k)=\frac{16\pi iN k e^{3dk}}{\alpha(2\cosh(2dk)-1)\sinh(6dk)} .
\]
The closest poles are double at \(\pm i\alpha\), with double-term coefficients \(N A_*, -N A_*\) there respectively, \(A_*=4\pi/(3\sqrt3 d^2)>0\). Shifting inverse contours past them gives terms \(-N A_* w e^{-\alpha w},-N A_* w e^{\alpha w}\) from the doubles on the right and left respectively, plus constant multiples of the respective exponentials times \(N\) and an exponentially faster remainder (uniform on the slightly wider strip since the transform decays horizontally there with room). Integrating to the endpoints and using principal logs in each edge coordinate \(u_e=N e^{\mp\alpha w}\) this says
\begin{equation}
      g-j_e=u_e\left[\kappa(\log N-\log u_e)+k_e+O((|u_e|/N)^\delta)\right]\qquad(|u_e|\le N),\qquad
      \kappa=A_*/\alpha^2>0                                                     \label{c4:Press-4}
\end{equation}
with constants \(\delta>0,k_e\) (the latter may differ by edge), valid even slightly beyond the right angles. For \(|u_e|\ge N\) on these sectors \(|g-j_e|\le C N\).

\paragraph{Projecting the logarithmic forcing.}
Keeping \(b_N\) as before, set
\[
              e_N=(f-g)/u_e,\qquad h_N=G/(u_e b_N),\qquad
              Q_N=(g-j_e)/(u_e b_N).
\]
The first two are analytic right with square-integrable boundary functions of positive Fourier-inverse support as above: for each \(N\), at zero \(f-g\) is \(O(|u_e|(1+|\log u_e|))\), \(G=O(|u_e|^{1/\alpha})\), and analogous inverse powers occur at infinity. Use \(P_+\) for the orthogonal projection on this boundary support \(t\ge0\). Then
\[
                       P_+(e_N/b_N)=-P_+ Q_N .
\]
Indeed the sum before projection is
\((f-j_e)/(u_e b_N)\), rational and analytic in the left half-plane,
regular at zero and decaying at infinity, so its Cauchy integral to the
right vanishes. Moreover in unweighted boundary norm of \(u_e\),
\[
 \Re\langle e_N,G/u_e\rangle\ge c_T\|e_N\|^2,\qquad
 \|h_N\|\le C\|e_N\|,\qquad \langle e_N,G/u_e\rangle=-\langle P_+Q_N,h_N\rangle .
\]
These follow by weighted \eqref{c4:Press-1} after dividing by \(u_e\) \((G=\mathsf T*(f-g))\), then by the boundary phase and support.

We check that \(P_+Q_N\) is uniformly bounded and converges weakly with right Laplace value at \(v'>0\) equal to
\[
                       \kappa\int_0^\infty \frac{e^{-v't}}{b+t}dt .
\]
Its Cauchy integral with denominator \(v'-u_e\), up the imaginary axis with factor \(1/(2\pi i)\), can be rotated to rays \(\arg u_e=\pm(\pi/2+\epsilon)\) for small fixed \(\epsilon>0\), oriented in and out respectively on lower and upper. There are no poles between and \eqref{c4:Press-4} justifies the endpoint passages. On those rays subtract \((\kappa\log N+k_e)e^{bu_e}\), which has zero integral by closing left. The resulting data before the Cauchy kernel are bounded uniformly in \(L^2\) there and converge weakly to \(-\kappa e^{bu_e}\log u_e\). Indeed for \(r=|u_e|\le N\) use \eqref{c4:Press-4} and
\[
        |b_N^{-1}|\le e^{-cr},\qquad |b_N^{-1}-e^{bu_e}|\le C(r^2/N)e^{-c' r}
\]
on the rays (the expansion of \(B_1(u_e/N)^{-1}\) near zero, with strict
modulus suppression along the left rays). For \(r\ge N\),
\(Q_N=O((N/r)e^{-cN^2/r})\), and
\[
 \int_N^\infty (N/r)^2e^{-2cN^2/r}\,dr
       =\int_0^N e^{-2c v}\,dv\le (2c)^{-1},\qquad v=N^2/r.
\]
Thus these tails are uniformly bounded in \(L^2\) and converge weakly
to zero: their support leaves every compact interval, so one first tests
against compactly supported \(L^2\) functions and then uses density.

The ray Cauchy evaluations are uniformly bounded in \(L^2\) up to the
imaginary boundary from the right, by angular separation: the kernel
bound \(C/(r+|\Im v'|)\) is bounded by the Schur test with power weight
\(-1/2\). They therefore bound the boundary projections by Plancherel,
and convergence of interior evaluations determines weak convergence.
Rotating the limiting rays to the two banks of the negative real axis
gives the log jump
\(\kappa\int_0^\infty e^{-b s'}/(v'+s')\,ds'\). By Fubini this is
\(\kappa\int_0^\infty e^{-v't}/(b+t)\,dt\), the claimed Laplace value.

\paragraph{The prescribed density tail and its sign.}
Consequently \(e_N,h_N\) are uniformly bounded in \(L^2\). Taking weak
subsequences as before, their limits give
\[
 \lim e_N(u_e)=\int_0^\infty e^{-u_e t}\varphi(t)dt,\qquad
 G_{\rm edge}(u_e)=\lim u_e b_N h_N=u_e\int_0^\infty e^{-u_e t}\psi(t)dt
\]
on the right with \(\varphi,\psi\in L^2(dt)\) and \(\psi=0\) for
\(t<b\). Dividing by the convergent boundary phases shifts the limiting
density left by \(b\). The projection identity and the just-computed
Laplace value therefore give
\[
                   \varphi(t)=-\kappa/t\qquad(t>b).
\]
This replaces the constant tail \(\varphi=1\) in the essential response.
Again \eqref{c4:Press-1} dilates to \(\psi=\mathsf T*\varphi\) in
\(x'=\log(t/b)/\alpha\); here the bound
\(|u_e e_N(u_e)|\le C\sqrt{u_e}\) on positive \(u_e\) justifies the
passage to the limit. The common intermediate function
\[
 q_{\rm edge}=\mathsf T_+^{-1}*\psi=\mathsf T_-*\varphi
   =\begin{cases}
       0,&x'<0,\\
       -(\kappa/b)c_-e^{-\alpha x'},&x'>0,
     \end{cases}
 \qquad c_-=\int e^{\alpha s'}\mathsf T_-(ds')>0
\]
is determined by the same two support arguments. It determines both edge
limits uniquely. Applying \(\mathsf T_+\) and integrating now gives
\[
        \int_0^\infty G_{\rm edge}(u_e)du_e=\int_b^\infty \psi(t)t^{-2}dt
             =-\frac{\kappa c_-}{2 b^2}\int e^{-\alpha s'}\mathsf T_+(ds')<0 .
\]
We recover \(J(id)=(L*S)(id)+(L*G)(id)\); the first term is \(N\gamma_0\) for a fixed constant. In the second, on real \(x\) we have \(|G(x)|\le C\sqrt{u_e}\exp(-c u_e)\) for \(u_e=N e^{-\alpha|x|}\) by the \(h_N\) bound. Thus exactly the same kernel limit as in \eqref{c4:Press-3} applies to that term, giving
\begin{equation}
                       J(id)=N\gamma_0-\gamma_1/N+o(1/N),\qquad \gamma_1>0 .  \label{c4:Press-5}
\end{equation}

\subsection{Positive activities and transverse-height estimates}\label{c4:transverse}

The physical response identities at the start of the section, now in
the \(w\)-coordinate, and \eqref{c4:Press-3} give
\begin{equation}
             A_N=J_1(id)\sim c/N\quad(c>0),\qquad
             J(id)=N(\log t_{\rm phys})'(d)+8 A_N^0,                          \label{c4:Press-6}
\end{equation}
where \(t_{\rm phys}(\lambda)\) is the empty-cell weight at \(\lambda,u=Q_\lambda^{-1}\).

In particular \(E_N=N^{-1}\sum_{\mathrm{ess}/v}|\gamma|w(\gamma)\ge A_N\), by translating marked edges to one representative row across the \(N\) columns, since each winding lift travels displacement \(\pm Nh\). So \eqref{c4:Press-3},\eqref{c4:Press-6} supply a polynomial lower bound for the marked identity before its full asymptotics. The same pressure also shows the plane SAW convergence radius cannot exceed \(x_*\). Indeed mark a \(P\)-vertex of a winding polygon, put its row at zero and lift the positively translating orientation to a plane SAW with starting column among \(N\) possibilities. Thus \(A_N\le N\sum_{l\ge N}c_l x_*^l\) for \(c_l\) the number of \(l\)-step plane SAWs from a given \(P\)-vertex. A strictly larger convergence radius would force exponential decay, contrary to \eqref{c4:Press-6}. This recovers the upper critical-radius bound, consistently with the connective-constant theorem cited above.

\begin{proposition}[Thin cylinders suppress large transverse spans]\label{prop:thin-cylinder-span}
Let $H(\gamma)$ be the span of cell row indices on the cylinder of period $Nh$. For every $\delta>0$, fixed $p\ge0$ and $K>0$, for both winding and contractible polygons,
\begin{equation}
                   \sum_{\gamma/v,\ H(\gamma)\ge N^{1+\delta}}
                         (1+|\gamma|)^p w(\gamma)=O(N^{-K}) .                 \label{c4:Press-7}
\end{equation}
\end{proposition}

\begin{proof}
\emph{Connected tuples.}
In the winding-only finite gas at \(\lambda=d\), the \(m\)-th log coefficient in \(s\) is \(1/m!\) times a weighted sum over ordered \(m\)-tuples (repetitions allowed), with additional factor the sum of \((-1)^{\#\mathrm{edges}}\) over connected spanning subgraphs of their incompatibility graph (edges indicate sharing a vertex). This follows by expanding pair-exclusion factors and exponentiating the connected diagrams. For a connected incompatibility graph the factor has sign \((-1)^{m-1}\) and magnitude at least 1, by deletion/contraction of a single edge and induction (parallel edges can be merged since the nonempty-subset sign sum is \(-1\)). For disconnected it is zero. Passing per height to axial translation classes of tuples (monotone in magnitude as above) and using \eqref{c4:Press-2} on the limiting log coefficient bounds their total connected-tuple weights by \(C_m/N\).

\emph{Winding spans.}
For an absolute winding \(\gamma\) let \(M_\gamma\) be total activity of absolute winding neighbors sharing a vertex with it, including itself. If \(H(\gamma)>2\), along a lift one can take a subpath between two same-sublattice vertices whose displacement \(D_\gamma=i'h+j'v\) has \(|j'|\ge H(\gamma)-2\) (adjust extremes by incident edges if necessary). For any other winding type \(\eta'\), a lifted period realizes \(Nh\). Periodize the two paths by their displacements; they intersect for every relative translation, since their parametrization difference is an invertible linear map on \(\mathbb R^2\) plus a bounded continuous periodic error and thus onto (e.g. by winding on a large circle). At lattice translates intersection gives shared vertices. Hence for every coset of the lattice modulo \(Nh,D_\gamma\) at least one translation of \(\eta'\) there meets the actual \(\gamma\) on the cylinder. There are \(N|j'|\) cosets. If the cylinder polygon \(\eta'\) has horizontal stabilizer of order \(k\), this gives at least \(N|j'|/k\) distinct neighbors in its full translation orbit; that orbit has \(N/k\) classes modulo axial translation. Vertical stabilizers are impossible. Thus \(M_\gamma\ge(H(\gamma)-2) A_N\). Restricting the tuple bound to stars gives
\[
         \sum_{\mathrm{ess}/v} w(\gamma)(1+H(\gamma))^{m-1}\le C'_m N^{m-2}
\]
by \eqref{c4:Press-6}. Since \(|\gamma|\le2N(H+1)\), this proves \eqref{c4:Press-7} for windings.

\emph{Contractible spans.}
For a tall contractible \(\delta'\), choose one \(R_{i,j}\) used in its highest row \(j\) (there is such an \(R\), and it uses both bonds within the row). Windings with lowest row \(j+1\) and occupying \(P_{i,j+1}\) have total absolute weight at least \(A_N/N\), by translating each axial class to that lowest row and averaging its nonempty lowest \(P\)-set by columns. Choose such a partner. Take symmetric edge difference of the polygon union with the hexagonal face whose lower and upper paths are respectively \(R_{i,j},P_{i+1,j},R_{i+1,j}\) and \(P_{i,j+1},R_{i,j+1},P_{i+1,j+1}\). Each polygon's overlap with the face consists of a single path of one or two edges including all occupied face vertices of that polygon (if the other extreme endpoint is occupied it also uses both within-row bonds). Thus this splices into a single simple winding \(\zeta\) (winding mod two is unchanged), with \(H(\zeta)\ge H(\delta')-1\), length within 2 of the sum. At most \(C|\zeta|\) inverse joins per output axial class are possible by choosing the touched face. Hence \(A_N/N\) times the tall contractible weight is bounded by the corresponding slightly relaxed tall essential sum with factor \(C|\zeta|\), and the length powers can be included on both sides. This proves \eqref{c4:Press-7}.

\end{proof}

Finally a plane translation class of simple polygons \emph{survives} period \(Nh\) when projection remains simple (otherwise there is a collision at lattice vertices). Exactly these project to cylinder contractibles, bijectively after full translations on both sides, by lifting. Each surviving plane class yields \(N\) cylinder axial classes, since a contractible lift is finite with disjoint period-translates, forbidding nontrivial horizontal stabilizers by lifting a coincidence. Therefore \(A_N^0/N\) is the unrooted plane activity sum over survivors. Each plane polygon eventually survives; \eqref{c4:Press-5},\eqref{c4:Press-6} and Fatou followed by bounded domination give finite total unrooted plane activity exactly \((\gamma_0-(\log t_{\rm phys})'(d))/8\). In particular the plane loss
\begin{equation}
                  \sum_{\mathrm{not~surviving}} w(\gamma)
                           \sim (\gamma_1/8) N^{-2}                           \label{c4:Press-8}
\end{equation}
over translation classes.
